\SourceLeaf{434}%
also
\[
\begin{aligned}
\frac{\partial X}{\partial x}&=Z^{3}\{-(1+uu)T+tuU\},&\quad\frac{\partial Y}{\partial x}&=Z^{3}\{tuT-(1+tt)U\},\\
\frac{\partial X}{\partial y}&=Z^{3}\{-(1+uu)U+tuV\},&\frac{\partial Y}{\partial y}&=Z^{3}\{tuU-(1+tt)V\}
\end{aligned}
\]
und
\[
\begin{aligned}
\omega&=Z^{6}(TV-UU)((1+tt)(1+uu)-ttuu)\\
&=Z^{6}(TV-UU)(1+tt+uu)\\
&=Z^{4}(TV-UU)\\
&=\frac{TV-UU}{(1+tt+uu)^{2}},
\end{aligned}
\]
genau derselbe Ausdruck, den wir am Schluss des vorigen Artikels gefunden haben. Man sieht also, dass
\begin{quote}
»das Krümmungsmass allemal durch den Bruch ausgedrückt wird, dessen Zähler $=1$, der Nenner das Product des grössten und kleinsten Krümmungshalbmessers in den durch die Normale gehenden Ebenen«.
\end{quote}

\Needspace{5\baselineskip}
\begin{center}12.\end{center}
\noindent\hspace*{1.2em}Wir wollen jetzt die Natur der kürzesten Linie auf der gekrümmten Fläche untersuchen. Die Natur einer gekrümmten Linie im Raum wird im Allgemeinen dadurch bestimmt, dass man die Koordinaten $x,y,z$ jedes Punkts als Funktionen einer einzigen veränderlichen Größe betrachtet, die wir $w$ nennen wollen. Die Länge der Kurve bis zu diesem Punkt, gemessen von einem beliebigen Anfangspunkt aus, ist dann
\[
=\int\BareRoot{\left\{\left(\frac{dx}{dw}\right)^{2}+\left(\frac{dy}{dw}\right)^{2}+\left(\frac{dz}{dw}\right)^{2}\right\}}.dw.
\]
Lässt man die Kurve ihre Lage um eine unendlich kleine Variation ändern, so ist die Variation der gesamten Länge
